\section{A uniform staircase at every integral power threshold}

\subsection*{Finding M-014: the mixed construction for all exponents at least six}
Let \(m\geq2\), \(c\geq6\) be integers and \(n\geq m^c\).
Every actual internally padded permanent multiplicity obstruction
\(q^m_\lambda(d[n])>K_n(\lambda)\) satisfies
\begin{equation}\label{eq:M014stairs}
             \lambda_{2j}\geq c+3-j
                         \qquad(1\leq j\leq c+1).
\end{equation}
Parts beyond the length are zero, as usual. This generalizes
the BIP staircase at \(c=25\). The separate IP staircase
at \(n>3m^4\) is already F-017 and is not obtained by
substituting \(c=4\) into the hypotheses above.

\begin{proof}
First, an actual obstruction requires \(m>2c\).
If \(4\leq m\leq2c\), then
\[
          m^c\geq4^c=2^{2c}\geq2^m>2^m-1.
\]
For \(m=2,3\), the same sufficient determinant size
\(2^m-1\) is at most \(m^c\) directly. The known
Grenet construction, in the checked form F-009, therefore
excludes every \(m\leq2c\). Also M-008 gives
\(r=\ell(\lambda)\geq2c+2\), while positive permanent
multiplicity gives \(r\leq m^2\).

We first prove the assertion at the cuts
\[
              k=2c+4-2s,\qquad 1\leq s\leq c+1,
                                                               \tag{M-014.1}
\]
except for \(s=c+1\) when \(c\) is odd. Suppose
\(\lambda_k\leq s\). Choose nonnegative integers
\(u,v\) with \(u+v=s-1\) and \(|u-v|\leq1\),
and put
\[
                      p=c-2u,\qquad q=c-2v.
\]
The stated range of \(s\), with its one exception,
ensures \(p,q\geq0\). Moreover \(p+q=k-2\).
Thus the F-017 mixed graph applies, with
\(R=r-k+1\), and has size
\[
 S=\binom{m+p}{p}\binom{R+u}{u}
      +\binom{m+q}{q}\binom{R+v}{v}-1.                 \tag{M-014.2}
\]
We show that each binomial product is strictly less than
\(m^c/2\).

Write
\[
 A_p(m)=\frac{\binom{m+p}{p}}{m^p},\qquad
 B_u(r)=\frac{\binom{r-1+u}{u}}{r^u}.
\]
The product formula shows that \(B_u(r)\) decreases
with \(r\). For \(r\geq2\),
\[
              B_u(r)\leq\frac{u+1}{2^u}\leq1.
                                                               \tag{M-014.3}
\]
For \(u\geq3\) and \(r\geq4\), the sharper bound is
\[
              B_u(r)\leq\frac{\binom{u+3}{3}}{4^u}
                                      \leq\frac5{16}.
                                                               \tag{M-014.4}
\]
The last sequence starts at \(20/64\) for \(u=3\)
and decreases, because its successive ratio is
\((u+4)/(4(u+1))<1\).

If \(p\geq3\), then \(m\geq2c\geq2p\) and
\[
 A_p(m)=\frac1{p!}\prod_{i=1}^p(1+i/m)
       \leq\frac{\exp(p(p+1)/(2m))}{p!}
       \leq\frac{\exp((p+1)/4)}{p!}
       \leq\frac e6<\frac12.
\]
The sequence in the penultimate expression decreases for
\(p\geq3\): its successive ratio is
\(e^{1/4}/(p+1)<1\). The elementary bound \(e<3\)
follows directly from its convergent factorial series.
Together with (M-014.3), this gives \(A_pB_u<1/2\).

If \(p=2\), the equation \(p+2u=c\geq6\) gives
\(u\geq2\). Since \(m\geq12\),
\[
 A_2(m)B_u(r)\leq\frac{91}{144}\frac34
                                      =\frac{91}{192}<\frac12.
\]
If \(p=1\), then \(u\geq3\), so (M-014.4) gives
\[
 A_1(m)B_u(r)\leq\frac{13}{12}\frac5{16}
                                      =\frac{65}{192}<\frac12.
\]
If \(p=0\), then \(u\geq3\) and the same bound
\(B_u(r)\leq5/16<1/2\) applies. These cases exhaust
the possibilities. Because \(R\leq r-1\), \(r\leq m^2\)
and \(p+2u=c\), they imply
\[
 \binom{m+p}{p}\binom{R+u}{u}
            \leq A_p(m)B_u(r)m^p r^u<\frac12m^c.
\]
The other product has the identical proof with \(q,v\).
Consequently \(S<m^c\leq n\). F-017 then gives
\(q^m_\lambda\leq K_n(\lambda)\), a contradiction.
Thus \(\lambda_k\geq s+1\) for every cut in
(M-014.1) covered so far.

The omitted odd-\(c\) endpoint is \(k=2\), \(s=c+1\).
It is precisely the jet consequence M-013:
an actual obstruction requires \(\lambda_2\geq c+2\).
All cases of (M-014.1) therefore hold. Substituting
\(j=c+2-s\) gives \eqref{eq:M014stairs}.
\end{proof}

\paragraph{Exact scope and Grothendieck dictionary.}
The result is a uniform numerical consequence of the actual
ideal-multiplicity comparisons F-017 and M-013, with the
known whole-permanent containment bound F-009 used only
to remove small \(m\). It does not infer multiplicity
inequalities from occurrence, nor does it impose the stronger-test
condition \(d>n/m\) on actual obstructions.
The torus-weight supports are the comodule gradings of
SGA~3, I, Proposition~4.7.3, pp.~35--36; their entire
highest-weight spaces arise from the Schur and exterior
operations of SGA~6, VI, Theorem~3.3, pp.~392--393.
The integral graph morphisms and the coefficient-jet map
retain the coordinate-algebra pullbacks. Vanishing of each
complete ideal comodule in characteristic zero is therefore
the precise scheme-theoretic content of the displayed
inequalities. No new purity or saturation input is used.
